\documentclass[12pt,a4paper]{article}
\usepackage{amsmath,amsthm,amssymb,amsfonts}
\usepackage[T1]{fontenc}
\usepackage[utf8]{inputenc}
\usepackage[ngerman]{babel}
\usepackage{hyperref}
\usepackage{geometry}
\geometry{margin=2.5cm}

\newtheorem{theorem}{Satz}[section]
\newtheorem{lemma}[theorem]{Lemma}
\newtheorem{proposition}[theorem]{Proposition}
\newtheorem{corollary}[theorem]{Korollar}
\newtheorem{definition}[theorem]{Definition}
\newtheorem{remark}[theorem]{Bemerkung}
\newtheorem{conjecture}[theorem]{Vermutung}

\begin{document}

\title{Die Collatz-Vermutung: Grundlagen, Iterationstheorie\\
und empirische Evidenz}
\author{Michael Fuhrmann}
\date{März 2026}
\maketitle

\begin{abstract}
Die Collatz-Vermutung, formuliert von Lothar Collatz im Jahr 1937, behauptet, dass jede
positive ganze Zahl unter wiederholter Anwendung einer einfachen arithmetischen Regel
schließlich die Zahl 1 erreicht. Trotz ihrer elementaren Formulierung bleibt die Vermutung
eines der bekanntesten offenen Probleme der Mathematik. Diese Arbeit stellt die grundlegende
Theorie der Collatz-Funktion $T\colon \mathbb{N} \to \mathbb{N}$ vor, untersucht ihr
Iterationsverhalten, diskutiert die Struktur von Zyklen, gibt einen Überblick über die
empirische Verifikation bis zu $2^{68}$ und beleuchtet Verbindungen zur Ergodentheorie,
Dichteargumenten und der Theorie der Vorgängermengen.
\end{abstract}

\tableofcontents

\newpage

\section{Einleitung}

Im Jahr 1937 stellte Lothar Collatz folgende verblüffend einfache Frage: Man nehme eine
beliebige positive ganze Zahl $n$. Ist sie gerade, teilt man durch 2. Ist sie ungerade,
multipliziert man mit 3 und addiert 1. Man wiederholt diesen Vorgang. Erreicht die
Folge immer schließlich die Zahl 1?

Diese Frage, heute bekannt als \emph{Collatz-Vermutung} (auch $3x+1$-Problem, Ulams
Vermutung oder Syrakus-Problem), hat fast neunzig Jahre lang allen Beweisversuchen
widerstanden. Paul Erd\H{o}s soll gesagt haben: ``Die Mathematik ist für solche Probleme
noch nicht bereit.''

Die Vermutung ist nicht nur wegen ihrer einfachen Formulierung bemerkenswert, sondern
auch wegen der Fülle mathematischer Strukturen, die sie angezogen hat: Zahlentheorie,
Ergodentheorie, $p$-adische Analysis, probabilistische Heuristiken und additive
Kombinatorik spielen alle eine Rolle in modernen Ansätzen.

\subsection{Notation und Konventionen}

In dieser Arbeit bezeichnet $\mathbb{N} = \{1, 2, 3, \ldots\}$ die positiven ganzen
Zahlen und $\mathbb{N}_0 = \{0, 1, 2, \ldots\}$ die nicht-negativen ganzen Zahlen.
Für eine Funktion $f\colon X \to X$ und $k \geq 0$ schreiben wir $f^k$ für das
$k$-fache Iterate von $f$, wobei $f^0 = \mathrm{id}$.

\section{Die Collatz-Funktion und ihre Iterationen}

\begin{definition}[Collatz-Funktion]\label{def:collatz}
Die \emph{Collatz-Funktion} (Standardform) ist die Abbildung $T\colon \mathbb{N} \to \mathbb{N}$,
definiert durch
\[
  T(n) \;=\; \begin{cases} n/2 & \text{falls } n \equiv 0 \pmod{2}, \\
                            3n+1 & \text{falls } n \equiv 1 \pmod{2}.
              \end{cases}
\]
Die \emph{Syrakus-Funktion} (oder komprimierte Form) fasst jeden ungeraden Schritt als
vollständigen Übergang von ungerade zu ungerade zusammen:
\[
  S(n) \;=\; \frac{3n+1}{2^{\nu_2(3n+1)}},
\]
wobei $\nu_2(m)$ die 2-adische Bewertung von $m$ bezeichnet (das größte $k$ mit $2^k \mid m$).
\end{definition}

\begin{remark}\label{rem:equivalent}
Die Vermutung ist in beiden Formulierungen äquivalent. Unter $S$ erscheinen nur ungerade
Zahlen, was für $p$-adische und Dichteargumente praktisch ist.
\end{remark}

\begin{definition}[Collatz-Folge und Stoppzeit]\label{def:sequence}
Für $n \in \mathbb{N}$ ist die \emph{Collatz-Folge} beginnend bei $n$ die Folge
$(n, T(n), T^2(n), \ldots)$.
Die \emph{totale Stoppzeit} $\sigma(n)$ ist das kleinste $k \geq 0$ mit $T^k(n) = 1$,
oder $\infty$, falls kein solches $k$ existiert.
Die \emph{Stoppzeit} $\sigma_\infty(n)$ ist das kleinste $k \geq 0$ mit $T^k(n) < n$,
oder $\infty$, falls kein solches $k$ existiert.%
\footnote{%
  \textbf{Nicht-Standard-Notation.} In dieser Arbeit bezeichnet $\sigma(n)$ die
  \emph{totale} Stoppzeit (erstes Iteriertes gleich 1) und $\sigma_\infty(n)$
  die \emph{partielle} Stoppzeit (erstes Iteriertes echt kleiner als $n$).
  Dies ist die \emph{umgekehrte} Konvention gegenüber Lagarias (1985) und
  Tao (2022), wo $\sigma_\infty(n)$ die totale Stoppzeit und $\sigma(n)$
  (bzw.\ $\sigma_T(n)$) die partielle Stoppzeit bezeichnet.
  Leserinnen und Leser, die jene Quellen konsultieren, mögen die beiden Symbole
  entsprechend vertauschen.%
}
\end{definition}

\begin{proposition}[Existenz von Vorgängern]\label{prop:pred}
Jede positive ganze Zahl $n$ hat mindestens einen Vorgänger unter $T$. Genauer gilt:
\begin{enumerate}
  \item[(i)] $2n$ ist stets ein direkter Vorgänger von $n$ (über den geraden Zweig): $T(2n) = n$.
  \item[(ii)] Falls $n \equiv 2 \pmod{3}$, so ist $m = (2n-1)/3$ eine positive ungerade
              ganze Zahl mit $T(m) = 2n$ und $T(2n) = n$, also ist $m$ ein Vorgänger
              von $n$ in zwei Schritten. Insbesondere hat $n$ einen ungeraden Vorgänger
              im Abstand 2 im Collatz-Graphen.
\end{enumerate}
\end{proposition}

\begin{proof}
(i) Per Definition gilt $T(2n) = 2n/2 = n$.\\
(ii) Setze $m = (2n-1)/3$. Dies ist genau dann eine positive ganze Zahl, wenn
$3 \mid (2n-1)$, d.h.\ $2n \equiv 1 \pmod{3}$, also $n \equiv 2 \pmod{3}$.
Da $3m = 2n-1$ ungerade ist, ist $m$ selbst ungerade. Es gilt
$T(m) = 3m + 1 = (2n-1) + 1 = 2n$ und $T(2n) = n$, also $T^2(m) = n$.
Somit ist $m$ ein Vorgänger von $n$ in zwei Schritten.\\
Hinweis: Ein direkter ungerader Vorgänger von $n$ über den ungeraden Zweig
würde $3m + 1 = n$ erfüllen, d.h.\ $m = (n-1)/3$; dies erfordert $n \equiv 1 \pmod{3}$
und $(n-1)/3$ ungerade.
\end{proof}

\section{Zyklen und Fixpunkte}

\begin{definition}[Zyklus]\label{def:cycle}
Ein \emph{Zyklus} von $T$ ist eine endliche Menge $\{n_0, n_1, \ldots, n_{k-1}\} \subset \mathbb{N}$
mit $T(n_i) = n_{i+1 \bmod k}$ für alle $i$. Die Zahl $k$ heißt die \emph{Periode}
des Zyklus.
\end{definition}

\begin{theorem}[Einziger bekannter Zyklus]\label{thm:cycle}
Der einzige in $\mathbb{N}$ bekannte Zyklus von $T$ ist der \emph{triviale Zyklus}:
Die Bahn $1 \to 4 \to 2 \to 1$ hat die Periode 3. In der Syrakus-Form entspricht dies
dem Fixpunkt $S(1) = 1$.
\end{theorem}

\begin{remark}\label{rem:cycle_search}
Alle Zyklen von $T$ mit Periode bis zu $10^{17}$ wurden rechnerisch untersucht
(vgl.\ \cite{Leavens1992}) und es wurde kein weiterer Zyklus außer dem trivialen gefunden.
\end{remark}

\begin{proposition}[Zyklusstrukturgleichung]\label{prop:cycle_eq}
Ein Zyklus ungerader Zahlen $m_0, \ldots, m_{a-1}$ unter $S$ mit 2-adischen Bewertungen
$e_i = \nu_2(3m_i + 1)$ und $E = e_0 + \cdots + e_{a-1}$ erfüllt
\begin{equation}\label{eq:cycle_cond}
  (2^E - 3^a)\, m_0 \;=\; C \;>\; 0,
\end{equation}
wobei $C > 0$ ein positiver ganzzahliger Korrekturterm ist, der von den Zwischenwerten abhängt.
Insbesondere muss für jeden Zyklus ungerader Zahlen unter $S$ die Ungleichung $2^E > 3^a$ gelten.
\end{proposition}

\begin{proof}
Unter $S$ gilt $m_{i+1} = (3m_i + 1) / 2^{e_i}$. Die $a$-fache Komposition liefert
\[
  m_0 = \frac{3^a m_0 + C}{2^E}
\]
für einen explizit berechenbaren positiven Term $C$, der von den Zwischenwerten abhängt.
Umstellen ergibt $(2^E - 3^a) m_0 = C > 0$, also $2^E > 3^a$.
Da jedoch $\log_2 3 \approx 1.585$ irrational ist, kann die Gleichung $2^E = 3^a$ nicht
exakt erfüllt werden. Die Ungleichung $2^E > 3^a$ erfordert $E/a > \log_2 3$, was
für kurze Zyklen ($a$ klein) strenge Bedingungen an die $e_i$ stellt und letztlich
erklärt, warum nicht-triviale Zyklen extrem selten oder inexistent sind.
\end{proof}

\section{Empirische Verifikation und Dichteargumente}

\begin{theorem}[Empirische Schranke]\label{thm:empirical}
Die Collatz-Vermutung wurde rechnerisch für alle positiven ganzen Zahlen bis zu
\[
  2^{68} \approx 2{,}95 \times 10^{20}
\]
verifiziert (vgl.\ \cite{Barina2020}). Das heißt: Für jedes $n \leq 2^{68}$ erreicht
die Collatz-Folge schließlich die Zahl 1.
\end{theorem}

\begin{remark}\label{rem:heuristics}
Heuristische Argumente stützen die Vermutung. Das ``durchschnittliche'' Verhalten von $T$
für eine ganze Zahl besteht in einem Multiplikationsfaktor von $3/2$ für ungerade Schritte
und $1/2$ für gerade Schritte. Da ungefähr die Hälfte aller Zahlen gerade ist, beträgt
der erwartete Faktor pro zwei Schritte ungefähr $(3/2)^{1/2} \cdot (1/2)^{1/2} = \sqrt{3/4} < 1$,
was geometrischen Abfall in Richtung 1 nahelegt.
\end{remark}

\begin{definition}[Logarithmische Dichte]\label{def:logdens}
Die \emph{logarithmische Dichte} einer Menge $A \subseteq \mathbb{N}$ ist
\[
  \mathrm{logDichte}(A) \;=\; \lim_{N \to \infty} \frac{1}{\log N} \sum_{\substack{n \leq N \\ n \in A}} \frac{1}{n},
\]
sofern der Grenzwert existiert. Dies ist ein schwächerer Dichtebegriff als die
natürliche Dichte $\lim_{N\to\infty} |A \cap [1,N]|/N$.
\end{definition}

\begin{proposition}[Positive Dichte gutartiger Bahnen]\label{prop:pos_density}
Für jedes feste $M \in \mathbb{N}$ hat die Menge
\[
  A_M = \{n \in \mathbb{N} : \exists k \geq 0,\; T^k(n) \leq M\}
\]
die natürliche Dichte 1. Das heißt: Fast alle positiven ganzen Zahlen (im Sinne der
natürlichen Dichte) steigen schließlich unter jede feste Schranke $M$ ab.
\end{proposition}

\begin{proof}
Wir geben zwei Argumente an: ein unbedingtes und ein bedingtes.

\textbf{Unbedingt} (nach Terras \cite{Terras1976}):
Man verwendet den Rückwärtsiterations-Baum mit Wurzel $M$. Für jedes $m \leq M$ wächst
die Menge der ganzen Zahlen, die innerhalb von $k$ Schritten $m$ erreichen können,
exponentiell in $k$ (mit Rate $\approx 2^k$), da jede positive ganze Zahl mindestens
einen Vorgänger unter $T$ besitzt (Proposition~\ref{prop:pred}).
Die komplementäre Menge der Zahlen, die innerhalb von $k$ Schritten \emph{nicht}
unter $M$ fallen, hat eine Dichte, die durch ergodische Abschätzungen
(siehe \cite{Terras1976}) durch $(3/4)^{k/2}$ nach oben beschränkt ist
(bis auf logarithmische Korrekturen). Im Grenzwert $k \to \infty$ ist die Dichte
der Ausnahmen gleich 0, also $\mathrm{Dichte}(A_M) = 1$.

\textbf{Bedingt} (unter Voraussetzung der Collatz-Vermutung):
Gilt die Collatz-Vermutung, so erreicht jede positive ganze Zahl schließlich 1;
damit liegt jede Zahl in $A_M$ für jedes $M \geq 1$. Die Menge
$\mathbb{N} \setminus A_M$ ist dann leer, und das Dichte-1-Resultat folgt unmittelbar.

\emph{Hinweis:} Das unbedingte Argument ist das mathematisch gehaltvolle; das bedingte
Argument zeigt lediglich, dass die Vermutung --- falls sie gilt --- ein wesentlich
stärkeres (punktweises) Ergebnis liefert.
\end{proof}

\section{Verbindungen zur Ergodentheorie}

\begin{definition}[Collatz-Dynamisches System]\label{def:shift}
Man betrachte die 2-adischen ganzen Zahlen $\mathbb{Z}_2$ (vgl.\ Abschnitt~\ref{sec:padic_preview}).
Die Collatz-Funktion setzt sich zu einer stetigen Abbildung
$\tilde{T}\colon \mathbb{Z}_2 \to \mathbb{Z}_2$
fort. Das dynamische System $(\mathbb{Z}_2, \tilde{T}, \mu)$, wobei $\mu$ das
Haar-Maß auf $\mathbb{Z}_2$ ist, liefert einen natürlichen ergodischen Rahmen.
\end{definition}

\begin{theorem}[Terras-Dichtesatz]\label{thm:terras}
Sei $N_k$ die Anzahl der positiven ganzen Zahlen $n \leq X$, für die $T^j(n) > n$ für
alle $j = 1, \ldots, k$ gilt. Dann
\[
  \frac{N_k}{X} \;\longrightarrow\; 0 \quad \text{für } X \to \infty,
\]
für jedes feste $k$. Genauer gilt $N_k/X \leq (3/4)^{k/2}$ bis auf logarithmische
Korrekturen.
\end{theorem}

\begin{proof}[Beweisansatz]
Die entscheidende Beobachtung ist, dass für eine ``zufällige'' ganze Zahl $n$ die
Stoppzeit $\sigma_\infty(n)$ unter einem geeigneten probabilistischen Modell mit
Wahrscheinlichkeit 1 endlich ist. Die Stoppzeit $\sigma_\infty(n)$ ist der erste
Zeitpunkt $k$, an dem der Zufallsgang $\sum_{j=1}^k X_j$ negativ wird, wobei
$X_j = -1$ mit Wahrscheinlichkeit $1/2$ (gerader Schritt) und
$X_j = \log_2(3/2)$ mit Wahrscheinlichkeit $1/2$ (ungerader Schritt).
Der Erwartungswert von $X_j$ ist $(1/2)\log_2(3/4) < 0$. Nach dem Gesetz der
großen Zahlen driftet der Gang nach $-\infty$, sodass er fast sicher die 0 unterschreitet.
Die Schranke $(3/4)^{k/2}$ folgt aus Abschätzungen großer Abweichungen.
\end{proof}

\begin{theorem}[Tao 2022 — fast beschränkte Bahnen]\label{thm:tao}
Für jede Funktion $f\colon \mathbb{N} \to \mathbb{R}_{>0}$ mit $f(n) \to \infty$ für
$n \to \infty$ hat die Menge
\[
  \bigl\{ n \in \mathbb{N} : \exists k \geq 0,\; T^k(n) \leq f(n) \bigr\}
\]
die logarithmische Dichte 1. Mit anderen Worten: Für fast alle $n$ (im Sinne der
logarithmischen Dichte) nimmt die Collatz-Bahn von $n$ fast beschränkte Werte an.
\end{theorem}

\begin{proof}[Beweisüberblick]
Taos Beweis \cite{Tao2022} (veröffentlicht in \emph{Forum of Mathematics, Sigma})
kombiniert Fourier-analytische Abschätzungen von
Exponentialsummen über die Syrakus-Abbildung mit Methoden der Ergodentheorie:
\begin{enumerate}
  \item Reduktion auf die Syrakus-Funktion $S$ auf ungeraden Zahlen.
  \item Modellierung der 2-adischen Bewertungen $e_i = \nu_2(3m_i+1)$ als annähernd
        unabhängige geometrische Zufallsvariablen.
  \item Einsatz der logarithmischen Fourier-Transformation zur Behandlung von
        Abhängigkeiten zwischen Schritten.
  \item Nachweis, dass das logarithmische Maß derjenigen $n$, die nicht unter
        $f(n)$ absinken, gleich 0 ist.
\end{enumerate}
Dies ist das stärkste bisher bewiesene unbedingte Resultat in Richtung der
Collatz-Vermutung.
\end{proof}

\section{Ausblick: $p$-adische und probabilistische Perspektiven}\label{sec:padic_preview}

\begin{definition}[2-adische ganze Zahlen]\label{def:2adic}
Die \emph{2-adischen ganzen Zahlen} $\mathbb{Z}_2$ sind die Vervollständigung von
$\mathbb{Z}$ bezüglich der 2-adischen Metrik $|x|_2 = 2^{-\nu_2(x)}$. Jedes Element
von $\mathbb{Z}_2$ lässt sich als formale Potenzreihe $\sum_{k=0}^\infty a_k 2^k$
mit $a_k \in \{0,1\}$ schreiben.
\end{definition}

\begin{proposition}[Fortsetzung von $T$ auf $\mathbb{Z}_2$]\label{prop:extension}
Die Collatz-Funktion $T\colon \mathbb{N} \to \mathbb{N}$ setzt sich zu einer stetigen
Abbildung $\tilde{T}\colon \mathbb{Z}_2 \to \mathbb{Z}_2$ fort. Die Unterscheidung
gerade/ungerade ist in $\mathbb{Z}_2$ durch das letzte Bit $a_0$ wohldefiniert.
\end{proposition}

\begin{proof}
Für $x \in \mathbb{Z}_2$ setze $\tilde{T}(x) = x/2$ falls $a_0 = 0$, und
$\tilde{T}(x) = (3x+1)/2$ falls $a_0 = 1$ (in der Syrakus-Form, die
$\mathbb{Z}_2^{\text{ungerade}}$ in sich abbildet). Die Abbildung ist stetig, da
beide Zweige 2-adisch stetig sind: Division durch 2 ist 2-adisch kontrahierend,
und $3x+1$ ist eine stetige affine Abbildung.
\end{proof}

\begin{conjecture}[Collatz-Vermutung]\label{conj:collatz}
Für jedes $n \in \mathbb{N}$ existiert ein $k \geq 0$ mit $T^k(n) = 1$.
Äquivalent dazu: Jede positive ganze Zahl liegt in der Rückwärtsbahn von 1 unter $T$.
\end{conjecture}

\section{Zusammenfassung und offene Fragen}

Die Collatz-Vermutung gehört zu den verlockendsten ungelösten Problemen der Mathematik.
Wir haben den grundlegenden Rahmen vorgestellt: die Collatz-Funktion $T$, ihre Iterationen,
die Struktur von Zyklen (eingeschränkt durch die Irrationalität von $\log_2 3$), empirische
Evidenz bis $2^{68}$, Dichteargumente für das gutartige Verhalten fast aller Bahnen und die
Verbindung zur Ergodentheorie, die in Taos Resultat von 2022 gipfelt.

Offene Fragen bleiben:
\begin{enumerate}
  \item Erreicht jede positive ganze Zahl schließlich die 1 (die Vermutung selbst)?
  \item Gibt es Zyklen außer dem trivialen $1 \to 4 \to 2 \to 1$?
  \item Kann Taos Resultat von logarithmischer Dichte auf natürliche Dichte verbessert werden?
  \item Wie sind die Stoppzeiten $\sigma(n)$ genau verteilt?
  \item Können $p$-adische Methoden eine Beweismethode liefern (vgl.\ \cite{Fuhrmann2026b})?
\end{enumerate}

\begin{thebibliography}{99}

\bibitem{Collatz1937}
L.~Collatz.
\newblock Unveröffentlichtes Problem, vorgestellt beim Internationalen Mathematikerkongress, 1937.

\bibitem{Lagarias1985}
J.~C. Lagarias.
\newblock The $3x+1$ problem and its generalizations.
\newblock \emph{American Mathematical Monthly}, 92(1):3--23, 1985.

\bibitem{Lagarias2010}
J.~C. Lagarias (Hrsg.).
\newblock \emph{The Ultimate Challenge: The $3x+1$ Problem}.
\newblock American Mathematical Society, Providence, RI, 2010.

\bibitem{Tao2022}
T.~Tao.
\newblock Almost all Collatz orbits attain almost bounded values.
\newblock \emph{Forum of Mathematics, Sigma}, 10:e72, 2022.

\bibitem{Terras1976}
R.~Terras.
\newblock A stopping time problem on the positive integers.
\newblock \emph{Acta Arithmetica}, 30(3):241--252, 1976.

\bibitem{Barina2020}
D.~Barina.
\newblock Convergence verification of the Collatz problem.
\newblock \emph{The Journal of Supercomputing}, 77:2681--2688, 2021.

\bibitem{Leavens1992}
G.~T. Leavens und M.~Vermeulen.
\newblock $3x+1$ search programs.
\newblock \emph{Computers \& Mathematics with Applications}, 24(11):79--99, 1992.

\bibitem{Furstenberg1977}
H.~Furstenberg.
\newblock Ergodic behavior of diagonal measures and a theorem of Szemer\'edi on arithmetic progressions.
\newblock \emph{Journal d'Analyse Math\'ematique}, 31(1):204--256, 1977.

\bibitem{Fuhrmann2026b}
M.~Fuhrmann.
\newblock Die Collatz-Vermutung und $p$-adische Zahlen.
\newblock Preprint, 2026.

\end{thebibliography}

\end{document}
